\RequirePackage[thinlines]{easybmat} %-- muss aufgrund eines Bugs vor etex und tikz geladen werden
\documentclass[a4paper, twoside, headsepline, index=totoc,toc=listof, fontsize=10pt, cleardoublepage=empty, headinclude, DIV=12, BCOR=5mm, titlepage]{scrartcl}
\usepackage{scrtime} % KOMA, Uhrzeit ermoeglicht
\usepackage{scrpage2} % wie fancyhdr, nur optimiert auf KOMA-Skript, leicht andere Syntax
\usepackage{etoolbox}
\usepackage{letltxmacro}
\usepackage{ifthen}

%-- Compiler-abhaengige IFs
\usepackage{ifxetex,ifluatex}
\newif\ifxetexorluatex
\ifxetex
  \xetexorluatextrue
\else
  \ifluatex
    \xetexorluatextrue
  \else
    \xetexorluatexfalse
  \fi
\fi

%--Farbdefinitionen
%-- muss vor tikz geladen werden
\usepackage[usenames, table, x11names]{xcolor}
\definecolor{dark_gray}{gray}{0.45}
\definecolor{light_gray}{gray}{0.6}

%--Zum Zeichnen
%-- muss vor polyglossia bzw. babel geladen werden
\usepackage{tikz}
\usepackage{tikz-cd}
\usetikzlibrary{external}
\tikzset{>=latex}
\usetikzlibrary{shapes,arrows,intersections}
\usetikzlibrary{calc,3d}
\usetikzlibrary{decorations.pathreplacing,decorations.markings, decorations.pathmorphing}
\usetikzlibrary{angles}

%-- Konfiguration von tikzexternalize
\tikzexternalize[prefix=tikz/, up to date check=diff]

\ifxetexorluatex
\tikzset{external/system call={xelatex \tikzexternalcheckshellescape %-- verwende LuaLaTeX, wegen dynamischer Speicherallokation
    -halt-on-error -interaction=batchmode -jobname "\image" "\texsource"}}
\else
\tikzset{external/system call={pdflatex \tikzexternalcheckshellescape 
    -halt-on-error -interaction=batchmode -jobname "\image" "\texsource"}}
\fi

%-- tikzexternalize fuer tikzcd deaktivieren, da inkompatibel
\AtBeginEnvironment{tikzcd}{\tikzexternaldisable}
\AtEndEnvironment{tikzcd}{\tikzexternalenable}

%-- um Inkompatibilitaeten von quotes und polyglossia bzw. babel zu vermeiden
\tikzset{
  every picture/.append style={
    execute at begin picture={\shorthandoff{"}},
    execute at end picture={\shorthandon{"}}
  }
}
\usetikzlibrary{quotes}
\usepackage{pgfplots}
\usepgfplotslibrary{colormaps}
\pgfplotsset{compat=1.10}



%-- Mathesymbole etc.
\usepackage{mathtools} % beinhaltet amsmath
\mathtoolsset{showonlyrefs, centercolon}
\newtagform{brackets}[\textbf]{[}{]}
\usetagform{brackets}
\usepackage{wasysym} % zusätzliche Symbole
\usepackage{amssymb} % zusätzliche Symbole
\usepackage{latexsym} % zusätzliche Symbole
\usepackage{stmaryrd} % für Blitz
\usepackage{nicefrac} % schräge Brüche
\usepackage{cancel} % Befehle zum Durchstreichen
\usepackage{mathdots} % Verbesserung von Punkten wie zB \ldots
\usepackage{mathrsfs} % das X wurde benötigt

%-- einzelne Symbole aus dem mathx-package
\DeclareFontFamily{U}{mathx}{\hyphenchar\font45}
\DeclareFontShape{U}{mathx}{m}{n}{<-> mathx10}{}
\DeclareSymbolFont{mathx}{U}{mathx}{m}{n}
\DeclareMathSymbol{\bigplus}{1}{mathx}{"90}

%-- 
\DeclareSymbolFont{bbold}{U}{bbold}{m}{n}
\DeclareSymbolFontAlphabet{\mathbbold}{bbold}
\newcommand{\mathds}[1]{\mathbb{#1}} % Um Kompatibilitaet mit frueheren Benutzung von dsfont herzustellen
\newcommand{\ind}{\mathbbold{1}} % charakteristische-Funktion-Eins

\def\mathul#1#2{\color{#1}\underline{{\color{black}#2}}\color{black}} %farbiges Untersteichen im Mathe-Modus
\renewcommand{\le}{\leqslant}
\renewcommand{\ge}{\geqslant}

%-- Underbrace als Befehl in LaTeX-Syntax (und ohne Spacing probleme mit nachfolgenden Operatoren...)
\newcommand{\Underbrace}[2]{{\underbrace{#1}_{#2}}}

%-- Alles was mit Schrift und XeTeX zu tun hat
\ifxetexorluatex
    % XeLaTeX or LuaTeX
	\ifxetex
		\usepackage{mathspec} %beinhaltet fontspec 
	\else
		\usepackage[no-math]{fontspec}
	\fi
    \usepackage{polyglossia} % moderner babel-ersatz
    \setmainlanguage[spelling=new,babelshorthands=true]{german}
	\shorthandoff{"}
	\setotherlanguage{english}
	% \newcommand\glqq{"}
	% \newcommand\grqq{"}
	\defaultfontfeatures{Mapping=tex-text, WordSpace={1.4}} %
    \setmainfont[Ligatures=Common, BoldFont={* Bold}, ItalicFont={* Light Italic},SmallCapsFont={Linux Libertine O}, SmallCapsFeatures={Letters=SmallCaps}]{Source Sans Pro}
	\setsansfont[Scale=MatchLowercase,Ligatures=Common, BoldFont={* Medium}]{Ubuntu}
	% \newfontface\fancy{texgyreadventor-italic.otf}
	% \setkomafont{title}{\fancy\Huge}
	\ifxetex
		\setallmonofonts[Scale=MatchLowercase, ItalicFont={*}]{Consolas} 
	\else
		\setmonofont[Scale=MatchLowercase, ItalicFont={*}]{Consolas}
	\fi
	\usepackage{xltxtra}
	\usepackage{fontawesome}
	\usepackage{microtype}
\else
    % default: pdfLaTeX
    \usepackage[ngerman]{babel}
    \usepackage[T1]{fontenc}
    \usepackage[utf8]{inputenc}
	\usepackage{textcomp} %verhindert ein paar Fehler bei den Fonts
    \usepackage[babel=true, tracking=true]{microtype}
	\usepackage{lmodern}
	\renewcommand{\familydefault}{\sfdefault} 
\fi


%--Mixed
\usepackage[neverdecrease]{paralist}
\usepackage[german=quotes]{csquotes}
\usepackage{makeidx}
\usepackage{booktabs}
\usepackage{wrapfig}
\usepackage{float}
\usepackage[margin=10pt, font=small, labelfont={sf, bf}, format=plain, indention=1em]{caption}
\captionsetup[wrapfigure]{name=Abb. }
\usepackage{stackrel}
\usepackage{multicol}

\flushbottom


%--Unterstreichung
\usepackage[normalem]{ulem}
\setlength{\ULdepth}{1.8pt}

%--Indexverarbeitung
\newcommand{\bet}[1]{\textbf{#1}}
\newcommand{\Index}[1]{\textbf{#1}\index{#1}}
\makeindex
\setindexpreamble{{\noindent \itshape Die \emph{Seitenzahlen} sind mit Hyperlinks zu den entsprechenden Seiten versehen, also anklickbar} \ifxetexorluatex \faHandUp\fi \par \bigskip}
\renewcommand{\indexpagestyle}{scrheadings}


%--Marginnote & todonotes
\deffootnote[1.5em]{1.5em}{1.5em}{\textsuperscript{\thefootnotemark}\ }
\usepackage[fulladjust]{marginnote}
\renewcommand*{\marginfont}{\color{dark_gray} \itshape \footnotesize}
\usepackage[textsize=small]{todonotes}
\usepackage{ragged2e}
\renewcommand*{\raggedleftmarginnote}{\RaggedLeft}
\renewcommand*{\raggedrightmarginnote}{\RaggedRight}

%--Todonotes mit tikz/externalize kompatibel machen
\LetLtxMacro{\oldtodo}{\todo}
\renewcommand{\todo}[2][]{\tikzexternaldisable\oldtodo[#1]{#2}\tikzexternalenable}
\LetLtxMacro{\oldmissingfigure}{\missingfigure}
\renewcommand{\missingfigure}[2][]{\tikzexternaldisable\oldmissingfigure[{#1}]{#2}\tikzexternalenable}

%--Konfiguration von Hyperref pdfstartview=FitH, 
\usepackage[hidelinks, pdfpagelabels,  bookmarksopen=true, bookmarksnumbered=true, linkcolor=black, urlcolor=SkyBlue2, plainpages=false,pagebackref, citecolor=black, hypertexnames=true, pdfauthor={Jannes Bantje}, pdfborderstyle={/S/U}, linkbordercolor=SkyBlue2, colorlinks=false]{hyperref}

\ifxetexorluatex
\newcommand{\hrefsym}[2]{\href{#1}{\texttt{#2\,\faExternalLink}}}
\newcommand{\hrefsymmail}[2]{\href{#1}{\texttt{\faEnvelope\,#2}}}
\renewcommand{\url}[1]{\hrefsym{#1}{\nolinkurl{#1}}}
\fi
\providecommand{\hrefsym}{\href}
\providecommand{\hrefsymmail}{\hrefsym}


% -- QR-Codes
\usepackage{qrcode} %-- hinter hyperref laden!

%--Römische Zahlen
\newcommand{\RM}[1]{\MakeUppercase{\romannumeral #1{}}}
\renewcommand{\thesection}{\Roman{section}}

%%--Abkürzungen etc.
\newcommand{\light}{\text{\Large $\lightning$}}

%-- Definitionen von weiteren Mathe-Befehlen
\DeclareMathOperator{\re}{Re} %Realteil
\DeclareMathOperator{\im}{Im} %Imaginaerteil
\DeclareMathOperator{\id}{id} %identische Abbildung
\DeclareMathOperator{\Hom}{Hom} 
\DeclareMathOperator{\Mat}{Mat}
\DeclareMathOperator{\Eig}{Eig}



\newcommand{\mathd}{\ensuremath{\mathrm{d}\mkern-1.0mu}}



%--Beweisende
\newcommand{\bewende}{\ifmmode \tag*{$\square$} \else \hfill $\square$ \fi}


%--Skalarprodukt
\DeclarePairedDelimiterX\sprod[2]{\langle}{\rangle}{#1\,\delimsize\vert\,#2}
\DeclarePairedDelimiterX\skal[2]{\langle}{\rangle}{#1\,,\,#2}

%--Betrag, Gaußklammer
\DeclarePairedDelimiter{\abs}{\lvert}{\rvert}
\DeclarePairedDelimiter{\floor}{\lfloor}{\rfloor}
\DeclarePairedDelimiter{\ceil}{\lceil}{\rceil}

%--Norm
\DeclarePairedDelimiter\doppelstrich{\Vert}{\Vert}
\newcommand{\norm}[2][\relax]{
\ifx#1\relax \ensuremath{\doppelstrich*{#2}}
\else \ensuremath{\doppelstrich*{#2}_{#1}}
\fi}


%--Umklammern
\DeclarePairedDelimiter\enbrace{(}{)}
\DeclarePairedDelimiter\benbrace{[}{]}
\DeclarePairedDelimiter\lenbrace{<}{>}
\newcommand{\ssbrace}[1]{{\scriptscriptstyle\enbrace{#1}}}

%--Mengen
\DeclarePairedDelimiterX\mengenA[1]{\lbrace}{\rbrace}{#1}
\DeclarePairedDelimiterX\mengenB[2]{\lbrace}{\rbrace}{#1\, \delimsize\vert \, #2}

\makeatletter
\newcommand{\set}[2][\relax]{
\ifx#1\relax \ensuremath{
\mengenA*{#2}}
\else \ensuremath{%
  \mengenB*{#1}{#2}}
\fi}
\makeatother

\DeclareRobustCommand{\minwidthbox}[2]{%
  \ifmmode
    \expandafter\mathmakebox
  \else
    \expandafter\makebox
  \fi
  [\ifdim#2<\width\width\else#2\fi]{#1}%
}


%--offen und abgeschlossen
\newcommand{\off}{\! \stackrel[\text{offen}]{}{\subset} \!}
\newcommand{\abg}{\! \stackrel[\text{abg.}]{}{\subset} \!}

%--Differential
\newcommand{\diff}[2]{\ensuremath{\frac{{\partial #1}}{{\partial #2}} }}
\newcommand{\diffd}[2]{\ensuremath{\frac{\mathd #1}{\mathd #2} }}

%--Diagonale Linien für Matrizen
\newcommand\tikzmark[1]{%
  \tikzexternaldisable\tikz[overlay,remember picture,baseline] \node [anchor=base] (#1) {};\tikzexternalenable}

\newcommand\MyLine[3][]{%
  \tikzexternaldisable\begin{tikzpicture}[overlay,remember picture]
    \draw[#1] (#2.south) -- (#3.north);
  \end{tikzpicture}\tikzexternalenable}
  
%--Jordankasten
\newcommand{\Jord}{ \tikzexternaldisable \begin{BMAT}{|ccc|}{|ccc|} %
				0 & & \\ %
				1\tikzmark{z} &  \diagdown & \\ %
				& \tikzmark{u}1 \MyLine[thick]{z}{u} &  0 %
			\end{BMAT} \tikzexternalenable}
			
%--Jordankasten mit Lambda
\newcommand{\JordLa}[1][\empty]{%
	\tikzexternaldisable\ifthenelse{\equal{#1}{\empty}}
		{\begin{BMAT}{|ccc|}{|ccc|} %
				\lambda & & \\ %
				1\tikzmark{z} &  \diagdown & \\ %
				& \tikzmark{u}1 \MyLine[thick]{z}{u} &  \lambda  %
			\end{BMAT}}
		{\begin{BMAT}{|ccc|}{|ccc|} %
				\lambda_{#1}  & & \\ %
				1\tikzmark{z} &  \diagdown & \\ %
				& \tikzmark{u}1 \MyLine[thick]{z}{u} &  \lambda_{#1}  %
			\end{BMAT}}\tikzexternalenable}

%--Inhaltsverzeichnis
\usepackage[tocindentauto]{tocstyle}
\usetocstyle{KOMAlike}	
\shorthandon{"}
		
%--Punkte (als letztes laden)
\usepackage{ellipsis}

\newcommand{\Semester}{SoSe 2015}
\newcommand{\fach}{Analysis 2}
\newcommand{\prof}{Prof.\ Ebert}

\providecommand{\verfasser}{Tim Keil}

%--Konfiguration von scrheadings
\setheadsepline{1pt}[\color{light_gray}]
\pagestyle{scrheadings}
\clearscrheadfoot

\providecommand{\shortFach}{\fach}
\lehead{ \includegraphics[height=0.6 cm,keepaspectratio]{../!config/Bilder/Logo_WWU_Muenster_light_gray.pdf}}
\rehead{\rule{0cm}{0.6cm}\footnotesize \sffamily \color{light_gray} \verfasser{} -- Mitschrift \shortFach}
\lohead{\rule{0cm}{0.6cm} \footnotesize \sffamily \color{light_gray} Stand: \today \; \thistime[:]}
\rohead{\includegraphics[height=0.6 cm,keepaspectratio]{../!config/Bilder/fb10logo_gray.pdf}}


\ofoot[{ \color{dark_gray} \LARGE \sffamily \thepage}]{{ \color{dark_gray} \LARGE \sffamily \thepage}} %hier wir auch der plain Stil bearbeitet!
\automark{section}
\ifoot{ \color{dark_gray} \small \leftmark}

%--Metadaten
\providecommand{\mail}{keil.menden@web.de}
\author{\verfasser}
\titlehead{\includegraphics[height=1.5cm, keepaspectratio]{../!config/Bilder/Logo_WWU_Muenster.pdf}%
\hfill \includegraphics[height=1.3cm, keepaspectratio]{../!config/Bilder/fb10logo.pdf}}
\title{Skript \fach}
\subtitle{Mitschrift der Vorlesung  \enquote{\fach} von \prof}



\begin{document}
\maketitle
\pagenumbering{Alph}
\section*{Vorwort --- Mitarbeit am Skript}
Dieses Dokument ist eine Mitschrift aus der Vorlesung \enquote{\fach, \Semester}, gelesen von \prof. Der Inhalt entspricht weitestgehend dem Tafelanschrieb. Für die
Korrektheit des Inhalts übernehme ich keinerlei Garantie! Für Bemerkungen und Korrekturen -- und seien es nur Rechtschreibfehler -- bin ich sehr dankbar. 
Korrekturen bitte durch persönliches Ansprechen oder per Mail an keil.menden@web.de.
%\begin{itemize}
%	\item Persönliches Ansprechen in der Uni oder Mails an keil.menden@web.de (gerne auch mit annotieren PDFs) \url{https://github.com/JaMeZ-B/latex-wwu}.
%	\item \emph{Direktes} Mitarbeiten am Skript: Den Quellcode poste ich auf GitHub (siehe oben), also stehen vielfältige Möglichkeiten der Zusammenarbeit zur Verfügung:
%	Zum Beispiel durch Kommentare am Code über die Website und die Kombination Fork + Pull Request. Wer sich verdient macht oder ein Skript zu einer Vorlesung, die 
%	ich nicht besuche, beisteuern will, dem gewähre ich gerne auch Schreibzugriff.
	
%	Beachten sollte man dabei, dass dazu ein Account bei \url{github.com} notwendig ist, der allerdings ohne Angabe von persönlichen Daten angelegt werden kann. 
%	Wer bei GitHub (bzw. dem zugrunde liegenden Open-Source-Programm \enquote{\texttt{git}}) -- verständlicherweise -- Hilfe beim Einstieg braucht, dem helfe ich gerne 
%	weiter. Es gibt aber auch zahlreiche empfehlenswerte Tutorials im Internet.\footnote{zB. \url{https://try.github.io/levels/1/challenges/1}, ist auf Englisch, aber dafür 
%	interaktives LearningByDoing}
%	\item \emph{Indirektes} Mitarbeiten: \TeX-Dateien per Mail verschicken. 
	
%	Dies ist nur dann sinnvoll, wenn man einen ganzen Abschnitt ändern möchte (zB. einen alternativen Beweis geben), da ich die Änderungen dann per Hand einbauen muss! Ich freue mich aber auch über solche Beiträge!
%\end{itemize}

\newpage

\tableofcontents
\cleardoubleoddemptypage
\pagenumbering{arabic}
\setcounter{page}{1}


\section{Normierte und metrische Räume} % (fold)
\label{sec:Normierte und metrische Räume}

\begin{centering}
	$\mathbb{R} \hspace{1cm} \rightarrow{} \hspace{1cm}  \mathbb{R}^n \hspace{1cm}  \rightarrow{} \hspace{1cm}  $Normierte Vektorräume$ \hspace{1cm} \rightarrow \hspace{1cm}  $Metrische Räume
\end{centering}

\subsection{Der $\mathbb{R}^n$} % (fold)
\label{sub:der_mathbb_r_n}

\[
	\mathbb{R}^n = \left\{ \left. \begin{pmatrix} x_1 \\ \vdots \\ x_n \\ \end{pmatrix}  \right| x_i \in \mathbb{R} \right\}
\]

Der $\mathbb{R}^n$ ist ein $\mathbb{R}$-Vektorraum und $\dim_{\mathbb{R}}(\mathbb{R}^n)=n$.

Sei
\[
	e_i = \begin{pmatrix} 0 \\ \vdots \\ 0 \\ 1 \\ 0 \\ \vdots \\ 0 \end{pmatrix} \leftarrow \text{$i$-ter Eintrag}
\]

\((e_1,\dots,e_n)\) ist eine Basis von $\mathbb{R}^n$ und heißt {\bfseries Standardbasis}. 

Ist \(x = \begin{pmatrix} x_1 \\ \vdots \\ x_n \end{pmatrix}\), so gilt \(x=\sum^{n}_{i=1}x_ie_i\)

\subsubsection{Definition} % (fold)
\label{ssub:definition}

Das Skalarprodukt auf $\mathbb{R}^n$ ist die Abbildung \[ \langle\, .\,,\,.\,\rangle : \mathbb{R}^n \times \, \mathbb{R}^n \to \mathbb{R}\] mit 

\[
 x = \left< \begin{pmatrix} x_1 \\ \vdots \\ x_n \end{pmatrix} , \begin{pmatrix} y_1 \\ \vdots \\ y_n \end{pmatrix} \right> := \sum_{i=1}^{n}x_iy_i \in \mathbb{R} 
 \]
 
es gilt:
\begin{description}
	\item[(S1)] (Bilinearität) $\forall \, x_0,x_1,y_0,y_1, \in \mathbb{R}^n \hspace{4pt} \forall \, a,b \in \mathbb{R}$
	 \begin{eqnarray}
	 \langle x_0+x_1\,,\, y_0+y_1\rangle & =& \langle x_0\,,\, y_0\rangle + \langle x_1\,,\,y_1\rangle + \langle x_1\,,\, y_0\rangle + \langle x_1\,,\, y_1\rangle \\
	 \langle ax_0\,,\, by_0\rangle &=& ab \langle x_0\,,\, y_0\rangle 
	 \end{eqnarray}
	 \item[(S2)](Symmetrie) $\forall \, x,y \in \mathbb{R}^n$ gilt $\skal{x}{y} = \skal{y}{x}$
	
	\item[(S3)](positive Definitheit) $\forall \, x \in \mathbb{R}$ gilt $\skal{x}{x} \geq 0$ und $\skal{x}{x}=0$ nur dann, wenn $x=0$
\end{description}
Desweiteren gilt $\skal{e_i}{e_j}=\delta_y := \begin{cases} 1 & i=j \\ 0 &i\neq j \end{cases}$ \hspace{1cm} (Kronecker-Symbol)

Ist \[
	x=\sum_{i=1}^{n}x_ie_i \in \mathbb{R}^n \qquad (x_i \in \mathbb{R})
	\]
so ist \begin{gather*}
	\skal{x}{e_j} = \skal{\sum_{i=1}^{n}x_ie_i}{e_j} \stackrel{(S1)}{=} \sum_{i=1}^{n}x_i \skal{e_i}{e_j} = \sum^{n}_{i=1}x_i \delta_{ij} = x_j \\
	\Rightarrow x = \sum^{n}_{j=1} \skal{x}{e_j}e_j
\end{gather*}

\subsubsection{Definition} % (fold)
\label{ssub:definition}
Sei $x \in \mathbb{R}^n$. Die $l^2$-Norm von $x$ ist die Zahl
\[
	\norm{x}_2:= \sqrt{\skal{x}{x}} 
\]
Schreibweise manchmal auch: $\norm{x}_2 = \norm{x}$
% subsubsection definition (end)

\subsection{Euklidsche und normierte Vektorräume} % (fold)
\label{sub:euklidsche_und_normierte_vektorraume}

\subsubsection{Definition} % (fold)
\label{ssub:definition}
Sei $ V$  ein $\mathbb{R}$-Vektorraum. Ein {\bfseries Skalarprodukt} auf $ V$  ist eine Abbildung $\skal{.}{.}:V \times V \to \mathbb{R}$, welche die Axiome S1,S2,S3 erfüllt. \\
Ein {\bfseries euklidscher Vektorraum} ist ein Paar $(V,\skal{}{})$. $ V$  ist ein $\mathbb{R}$-Vektorraum und $\skal{}{}$ ist Skalarprodukt auf $ V$. \\
Ist $(V,\skal{}{})$ ein eukdlischer Vektorraum, so ist die durch $\skal{}{}$ induzierte Norm definiert als \[
	\norm{v}=\sqrt{\skal{v}{v}} \hspace{1cm}  v \in V 
	\]

\subsubsection{Beispiele} % (fold)
\label{ssub:beispiele}

\begin{itemize}
	\item $\mathbb{R}^n$ mit Standardskalarprodukt
	\item Die Abbildung \begin{align*}
		 \mathbb{C} \times \mathbb{C} &\to \mathbb{R} \\
		 (z,w) &\mapsto Re(\bar z \cdot w)
		 \end{align*}
		 ist ein Skalarprodukt. Die induzierte Norm ist der Absolutbetrag.      
\end{itemize}
% subsubsection beispiele (end)

\subsubsection{Definition} % (fold)
\label{ssub:definition}
Sei $ V$ ein $\mathbb{R} $-Vektorraum. Eine {\bfseries Norm} auf $ V$ ist eine Abbildung $\norm{.}:V \to \mathbb{R} , v \mapsto \norm{v}$ mit
\begin{description}
	\item[(N1)]$\norm{v} \geq 0 \hspace{1cm} \forall \, v \in V$ 
	\item[(N2)] $\norm{v}=0 \Leftrightarrow v=0$
	\item[(N3)] $\norm{av}=|a|\cdot \norm{v} \hspace{1cm}  \forall\,v \in V, a \in \mathbb{R}$
	\item[(N4)] $\norm{v+w} \leq \norm{v}+\norm{w} \hspace{1cm} \forall\, v,w \in V$ \hspace{1cm} (Dreiecksungleichung) 
\end{description} 
Ein normierter Vektorraum ist ein Paar $(V,\norm{.})$ bestehend aus einem $\mathbb{R}$-VR $ V$ und einer Norm auf $ V$.

\subsubsection{Satz} % (fold)
\label{ssub:satz}
Sei $(V,\skal{}{})$ ein euklidscher Vektorraum. Für $v \in V$ setzte $\norm{v}=\sqrt{\skal{v}{v}}$ \\ Dann gilt:
\begin{description}
	\item[1.)] $\norm{.}$ ist eine Norm auf $ V$ 
	\item[2.)] $\forall\, v,w \in V$ gilt die {\bfseries Cauchy-Schwarz-Ungleichung}:
	\[
		|\skal{v}{w}| \leq \norm{v} \cdot \norm{w}
	\]  
\end{description}
\minisec{Beweis:}
\begin{enumerate}
	\item \begin{description}
		\item[N1:] $\norm{v}=\sqrt{\underset{\geq 0}{\underbrace{\skal{v}{v}}}} \geq  0$
		\item[N2:] $\norm{v}=\sqrt{\skal{v}{v}}=0 \Rightarrow \skal{v}{v}=0 \stackrel{\text{S3}}{\Rightarrow} v=0$ 
		\item[N3:] $a \in \mathbb{R} , v \in V$ \[
			\norm{av}=\sqrt{\skal{av}{av}} \stackrel{\text{S1}}{=} \sqrt{a^2 \skal{v}{v}} = \sqrt{a^2} \sqrt{\skal{v}{v}} = \abs{a} \norm{v}
		\] 
	\end{description}
	\item CS-Ungleichung ist trivial, wenn $w=0$ \\
	Ann.: $w \neq 0$: Sei $f:\mathbb{R} \to \mathbb{R} $ mit
	\begin{align*}
		f(t) &=\skal{v+tw}{v+tw} \\
		& = \skal{v}{v} +2t \skal{v}{w} + t^2 \skal{w}{w} \\
		& = \norm{v}^2+ 2t \underbrace{\skal{v}{w}}_{\in \mathbb{R}} + t^2 \norm{w}^2 
	\end{align*} 
	aus S3 folgt $f(t)\geq 0$. f ist Polynom
	\begin{gather*}
		f'(t)=0 \Leftrightarrow t_0 = - \frac{\skal{v}{w}}{\norm{w}^2} \\
		f''(t)=2 \norm{w}^2 > 0
	\end{gather*}
 $\Rightarrow$ $f$ nimmt bei $t_0$ globales Minimum an
 \begin{gather*}
 	0 \leq f(t_0)=\norm{v}^2 - \frac{2 \skal{v}{w}^2}{\norm{w}^2} + \frac{\skal{v}{w}^2}{\norm{w}^4} \norm{w}^2 = \norm{v}^2 - \frac{\skal{v}{w}^2}{\norm{w}^2} \\
	\Rightarrow \frac{\skal{v}{w}^2}{\norm{w}^2} \leq \norm{v}^2 \\ \Rightarrow \skal{v}{w}^2 \leq \norm{v}^2 \norm{w}^2 \\
	\Rightarrow \abs{\skal{v}{w}} = \sqrt{\skal{v}{w}^2} \leq \sqrt{\norm{v}^2 \norm{w}^2} = \norm{v} \cdot \norm{w}
 \end{gather*}
 \bewende
\end{enumerate}
\begin{description}
	\item[(N4)]\begin{align*}
		\norm{v+w}^2 &= \skal{v+w}{v+w} \\ 
		&\stackrel{\text{S1}}{=} \norm{v}^2+ 2 \skal{v}{w}+ \norm{w}^2 \\
		&\leq \norm{v} + 2 \abs{\skal{v}{w}}+ \norm{w}^2 \\
		&\stackrel{\text{CS}}{\leq} \norm{v}^2+2 \norm{v}\norm{w}+ \norm{w}^2 \\
		&=(\norm{v}+\norm{w})^2
	\end{align*} 
Wurzel ziehen zeigt N4
\bewende
\end{description}
% subsubsection satz (end)

\subsubsection{Weitere Beispiele für Normen auf $\mathbb{R}$} % (fold)
\label{ssub:weitere_beispiele_fur_normen_auf_mathbb_r}

\begin{gather*}
	 x = \sum^{n}_{i=1}x_ie_i \hspace{1cm}  x_i \in  \mathbb{R}
	 \end{gather*}
\begin{itemize}
	\item \text{$l^1$-Norm} $\hspace{1cm}  \norm{x}_1 = \sum^{n}_{i=1}\abs{x_i}$ 
	\item $l^{\infty}$-Norm $\hspace{1cm} \norm{x}_{\infty}=\max{\{\abs{x_i}\,|\,i=1,\dots,n}\}$
	\item Ist $p \in [1,\infty)$, so ist die $l^p$-Norm
	\[
		\norm{x}_p = \left( \sum^{n}_{i=1}\abs{x_i}^p \right)^{\frac{1}{p}} \hspace{1cm} \text{(Beweis Schwer)}
	\] 
\end{itemize}
% subsubsection weitere_beispiele_fur_normen_auf_mathbb_r (end)

\subsubsection{Definition} % (fold)
\label{ssub:definition}
$ V$ $\mathbb{R}$-VR , $\norm{.}_0 , \norm{.}_1$ Normen auf $ V$. Die Normen $\norm{.}_0 , \norm{.}_1$ heißen {\bfseries äquivalent}, falls $C,c>0$ existieren, so dass für alle $v \in  V$ gilt:
\[
	c \norm{v}_0 \leq \norm{v}_1 \leq C \norm{v}_0 
\] 
\minisec{Bemerkung:} Das ist eine Äquivalenzrelation
\begin{itemize}
	\item {\bfseries Reflexiv}: klar
	\item {\bfseries Symmetrie}: \begin{gather*}
		c \norm{v}_0 \leq \norm{v}_1 \leq C \norm{v}_0 \\
		\Rightarrow \frac{1}{c} \norm{v}_1 \leq \norm{v}_0 \leq \frac{1}{C} \norm{v}_1
	\end{gather*}
	\item {\bfseries Transitiv}: \begin{gather*}
		c \norm{v}_0 \leq \norm{v}_1 \leq C \norm{v}_0 \hspace{0.5cm} \text{und}\hspace{0.5cm} d \norm{v}_1 \leq \norm{v}_2 \leq D \norm{v}_1  \hspace{1cm} (C,c,D,d > 0)	\end{gather*}
		\[
		\Rightarrow cd \norm{v}_0 \leq CD \norm{v}_0
		\]
\end{itemize}

\subsubsection{Beispiel} % (fold)
\label{ssub:beispiel}
Die $l^1, l^2 \, \text{und} \, l^{\infty}$-Norm auf $\mathbb{R}^n $ sind äquivalent und zwar gilt:
\[
	\norm{v}_{\infty} \leq \norm{v}_2 \leq \norm{v}_1 \leq n \norm{v}_{\infty} \hspace{1cm} \text{(Beweis Übung)}
\]
% subsubsection beispiel (end)

\subsection{Metrische Räume} % (fold)
\label{sub:metrische_raume}

\subsubsection{Definition} % (fold)
\label{ssub:definition}
$(V,\norm{.})$ normierter Vektorraum, $v,w \in  V$ \\
Der {\bfseries Abstand} von $v$ und $w$ ist \[
d(v,w):= \norm{v-w}	
\] 
% subsubsection definition (end)
 \subsubsection{Definition} % (fold)
 \label{ssub:definition}
Sei $X$ eine Menge \\
Eine {\bfseries Metrik} (oder Abstandsfunktion) ist eine Abbildung $d: X \times X \to \mathbb{R} , (x,y) \mapsto d(x,y)$, sodass gilt:
\begin{description}
	\item[(M1)] $\forall\, x,y \in X: \hspace{1cm}  d(x,y) \geq  0$
	\item[(M2)] $\forall\, x,y \in X: \hspace{1cm}  d(x,x)=0   \hspace{1cm} \text{und} \hspace{1cm} 
		 		d(x,y)=0 \Rightarrow x=y$
	\item[(M3)] $\forall\, x,y \in X: \hspace{1cm} d(x,y)=d(y,x)$
	\item[(M4)] $\forall\, x,y,z \in X: \hspace{1cm}  d(x,z) \leq d(x,y)+d(y,z) \hspace{1cm}$ ({\bfseries Dreiecksungleichung}) 
\end{description}
Ein {\bfseries Metrischer Raum} ist ein Paar $(X,d)$, $X$ Menge, $d$ Metrik auf $X$
\minisec{Beispiele:}
\begin{itemize}
	\item Ist $(V,\norm{.})$ normierter Vektorraum, so ist durch $d(v,w):= \norm{v-w}$ eine Metrik auf $ V$ gegeben.

\begin{align*}
	v,w,u \in V \hspace{1cm} d(v,u) &= \norm{v-u} = \norm{v-w+w-u} \\
	&\stackrel{\text{$\triangle$-Ungl.}}{\leq} \norm{v-w} + \norm{w-u} \\ 
	&= d(v,w)+d(w,u)
\end{align*}
 
\minisec{Unterbeispiel:}
$V=\mathbb{R}$. Dann ist $x \mapsto \abs{x} $ eine Norm auf $\mathbb{R}$ und die induzierte Metrik auf $\mathbb{R}$ ist $d(x,y)=\norm{x-y}$
 % subsubsection definition (end)
 \item $(X,d_X)$ metrischer Raum $Y\subseteq X$ Teilmenge. Dann ist $(Y,d_y)$ ein metrischer Raum, wobei für $y_0,y_1 \in Y \text{gilt, dass} \, d_y(y_0,y_1):=d_x(y_0,y_1)$ \\
 \underline{Vereinbarung:} Ist $X \subseteq \mathbb{R}^n$ eine Teilmenge, so ist $X$ immer mit der durch das Standardskalarprodukt induzierte Metrik versehen. (es sei denn, etwas anderes ist gesagt) 
\end{itemize}

\minisec{Konvergenz}
Erinnerung (Ana I): $(x_n)_n$ Folge in $\mathbb{R}$, $x \in  \mathbb{R}$. Dann 
\[
\lim_{n \to \infty}x_n=x \, :\Leftrightarrow \, \forall\, \varepsilon > 0 \,\exists \,n_0 \, \forall\, n \geq n_0: \, \abs{x-x_1}< \varepsilon
 \] 
% subsection metrische_raume (end)
\subsubsection{Definition} % (fold)
\label{ssub:definition}
$(X,d)$ metrischer Raum, $(x_n)_n$ Folge in $X$ (d.h eine Abbildung $x: \mathbb{N} \to X$) $y \in X$. Dann Konvergiert die Folge $(x_n)_n$ gegen $y :\Leftrightarrow \forall\, \varepsilon >0 \, \exists \, n_0: \forall \ n \geq n_0 \, d(x_n,y)< \varepsilon $
\underline{Äquivalent}: Die Folge (in $\mathbb{R}$) $d(x_n,y)$ ist eine Nullfolge.
% subsubsection definition (end)

\subsubsection{Beispiel} % (fold)
\label{ssub:3.4}
$X=\mathbb{R}^n$ mit Metrik, die durch $l^1$-$ l^2$-$ \text{ oder } l^{\infty}$- Norm induziert wird. Folge $k \mapsto x(k) \in \mathbb{R}^n$.
\[
	x(k)= \sum_{i=1}^{n}x_i(k)e_i \qquad e_i \text{ Standard-Basisvektor} \qquad x_i(k) \in  \mathbb{R} \qquad x=\sum^{n}_{i=1}x_ie_i
\]
Dann sind äquivalent
\begin{description}
	\item[1)] für alle $i=1,\dots,n$ gilt $\lim_{k \to \infty}x_i(k)=x_i$
	\item[2)] $\lim_{k \to \infty}\norm{x(k)-x}_1=0$
	\item[3)] $\lim_{k \to \infty}\norm{x(k)-x}_2=0$
	\item[4)] $\lim_{k \to \infty}\norm{x(k)-x}_{\infty}=0$ 
\end{description}
\minisec{Beweis:}
\underline{$1 \to 2$:} 
\[
	\norm{x(k)-x}_1:=\sum^{n}_{i=1}\underset{\stackrel{\text{(wg 1)}}{\to}0}{\underbrace{\abs{x_i(k)-x_i}}} \longrightarrow 0
\]
\\
\underline{$2 \to 3 \to 4$:} Benutze Ungleichungen $(y \in \mathbb{R}^n)$
\[
	\norm{y}_{\infty} \leq \norm{y}_2 \leq \norm{y}_1 \leq n \cdot \norm{y}_{\infty}
\]
also 
\begin{gather*}
	\lim_{k \to \infty}\norm{x(k)-x}_1=0 \\
	\Rightarrow \lim_{k \to \infty}\norm{x(k)-x}_2=0 \\
	\Rightarrow \lim_{k \to \infty}\norm{x(k)-x}_{\infty}=0
\end{gather*}
\\
\underline{$4 \to 1$:} \begin{gather*}
	\norm{x(k)-x}_{\infty} = \max\{ \abs{x_i(k)-x} \,| i=1,\dots,n \}
	\Rightarrow \text{für alle $i$: } \abs{x_i(k)-x_i} \leq \norm{x(k)-x}_{\infty}
\end{gather*}
Also
\[
	\left[ \norm{x(k)-x}_{\infty} \stackrel{k \to \infty}{\longrightarrow}0 \right] \qquad \Rightarrow \qquad \left[ \forall\, i=1,\dots,n \quad \abs{x_i(k)-x_i} \longrightarrow 0 \right]
\]
\[
\Rightarrow 1)
\] \bewende

\minisec{Bemerkung:} $X$ metrischer Raum $(x_n)_n , (y_n)_n$ Folgen in $X \quad \lim_{n \to \infty}x_n=x \quad \text{und} \quad \lim_{n \to \infty}y_n=y$
\[
	\Rightarrow d(x,y)=\lim_{n \to \infty}d(x_n,y_n) \qquad \text{(Blatt 2, Aufgabe 2)}
\]
Daraus folgt: \\ Ist $(x_n)_n $ Folge in $X$ und gilt $\lim_{n \to \infty}x_n=x$ sowie $\lim_{n \to \infty}x_n=x'$ dann ist $x=x'$ \\
Denn \[
	d(x,x')=\lim_{n \to \infty}\underset{=0}{\underbrace{d(x_n,x_n)}}=0
\]
\[
	\stackrel{(M2)}{\Rightarrow} x=x'
\] \bewende
% subsubsection beispiel (end)

\subsubsection{Definition} % (fold)
\label{ssub:definition}
Sei $(X,d)$ metrischer Raum und $(x_n)_n $ Folge in $X$. Dann heißt $(x_n)_n$ {\bfseries Cauchyfolge}, falls gilt:
\[
	\forall\, \varepsilon >0 \, \exists \, n_0 \, \forall\, n,m \geq n_0 \, d(x_n,x_m)<\varepsilon 
\]

\minisec{Bemerkung:} $(X,d)$ metrischer Raum $(x_n)_n $ Folgen in $X$, $\lim_{n \to \infty}x_n=x \in X$. \\
Dann ist $(x_n)_n$ Cauchyfolge, denn \\
Sei $\varepsilon >0$. Wähle $n_0$, sodass \[  \forall\, n \geq n_0 \, d(x,x_n)< \frac{\varepsilon}{2} \]
Falls $n,m \geq n_0$
\[
	d(x_n,y_n) \leq \underset{<\frac{\varepsilon}{2}}{\underbrace{d(x_n,x)}} + \underset{<\frac{\varepsilon}{2}}{\underbrace{d(x,y_n)}} < \varepsilon 
\]
% subsubsection definition (end)

\subsubsection{Definition} % (fold)
\label{ssub:definition}
$(X,d)$ metrischer Raum. Dann heißt $X$ {\bfseries vollständig}, falls jede Cauchy-Folge $(x_n)_n $ in $X$ gegen ein $x \in X$ konvergiert (Ana I)

\minisec{Bsp:} 

	$(0,1) \subseteq \mathbb{R}$.
	$x_n = \frac{1}{n}$ Folge in $X$ ist Cauchyfolge. \\
	$(x_n)_n$ konvergiert \underline{nicht} in $X$, denn der einzig möglicher Grenzwert (0) liegt nicht in $X$.

% subsubsection definition (end)

\subsubsection{Satz} % (fold)
\label{ssub:satz}
$\mathbb{R}^n$, mit der durch die Norm $\norm{.}_p \qquad p=1,2,\infty$ induzierten Metrik ist vollständig

\minisec{Beweis:}
\underline{$p=\infty$:} Sei $x(k)$ $l^{\infty}$-Cauchyfolge in $\mathbb{R}^n$. Dann gilt für $i=1,\dots,n$:
\[
	\abs{x_i(k)-x_i(l)} \leq \norm{x(k)-x(l)}_{\infty}
\]
Da $(x(k))_k$ $l^{\infty}$-Cauchyfolge, ist $(x_i(k))_k$ eine Cauchyfolge in $\mathbb{R}$ für alle $1\leq i \leq n$ \\
Nach dem Cauchy-Kriterium aus Analysis I sind die Folgen $(x_i(k))_k$ Konvergent. $x_i=\lim_{k \to \infty}x_i(k) \in \mathbb{R}$
\[
	x := \sum_{i=1}^{n}x_ie_i \in \mathbb{R}^n
\] 
Dann gilt 
\[
	\norm{x(k)-x}_{\infty} = \max\{ \underset{\to 0}{\underbrace{\abs{x_i(k)-x_i}}} | i=1,\dots,n \} \longrightarrow 0
\]
also
\[
	\lim_{k \to \infty}\norm{x(k)-x}_{\infty} = 0
\]
\\
\underline{$p=2$:}
$x(k)$ $l^2$-Cauchyfolge, weil $\norm{x(k)-x(l)}_{\infty} \leq \norm{x(k)-x(l)}_2$ \\
Ist $(x_k)_k$ eine $l^{\infty}$-Cauchfolge, so gibt es ein $x \in \mathbb{R}^n$ mit 
\[
	\lim_{k \to \infty}\norm{x(k)-x}_{\infty}=0
\]
\[
	\Rightarrow \norm{x(k)-x}_2 \leq n \cdot \underset{\to 0}{\underbrace{\norm{x(k)-x}_{\infty}}} \longrightarrow 0
\] 
\\
\underline{$p=1$:} Analog zum Fall $p=1$, benutze $\norm{y}_{\infty} \leq \norm{y}_1 \leq n \cdot \norm{y}_{\infty}$ 
\bewende	
% subsubsection satz (end)
\subsubsection{Definition} % (fold)
\label{ssub:definition}
Seien $(X,d_X)$ und $(Y,d_Y)$ metrische Räume und $f:X \to Y$ eine Abbildung sowie $x \in X$.
\begin{enumerate}[(i)]
	\item $f$ heißt {\bfseries $\varepsilon$-$\delta$-stetig} (oder {\bfseries stetig}) bei $x$, wenn gilt
	\[
		\forall\, \varepsilon >0 \, \exists \, \delta>0 \, \forall\, x' \in X \text{ mit } d_X(x,x')< \delta \text{ gilt: } d_Y(f(x),f(x'))<\varepsilon 
	\]
	\item $f$ heißt {\bfseries folgenstetig} in $x$ $:\Leftrightarrow$ Ist $(x_n)_n $ eine Folge in $X$ mit $\lim_{n \to \infty}x_n=x$, so gilt
	\[
		\lim_{n \to \infty}f(x_n)=f(x)
	\]
\end{enumerate}
$f$ ist stetig (folgenstetig) auf $X$, wenn $f$ stetig (folgenstetig) bei $x$ ist für alle $x \in X$.
% subsubsection definition (end)

\subsubsection{Satz} % (fold)
\label{ssub:satz}
$f$ stetig in x $\Leftrightarrow$ $f$ folgenstetig in $x$
\minisec{Beweis} wie in Analysis I, wird bald nachgeliefert

% subsubsection satz (end)
\subsubsection{Satz} % (fold)
\label{ssub:satz}
$(X,d_X)$ , $(Y,d_Y)$, $(Z,d_Z)$ metrische Räume. 
\[ f:X \to Y \qquad g:Y \to Z \qquad x \in X \]
Falls $f$ stetig bei $x$ und stetig bei $f(x)$ so ist $g \circ f$ stetig bei $x$.
\minisec{Beweis:} Siehe Analysis I 
% subsubsection satz (end)
% subsection euklidsche_und_normierte_vektorraume (end)

\end{document}	